\subsection{GameWorld - Die Spielwelt als Ressourcenverwalter}

Der \texttt{GameWorld} kommt in helios die Rolle als Datencontainer zu und besitzt als zentraler Hub den \texttt{EntityManager} sowie verschiedene Registries für den Zugriff auf \texttt{Manager}, \texttt{CommandHandler} und \texttt{CommandBuffer}.

Da sie Zugriff auf die relevanten Systeme ermöglicht, ähnelt sie in Teilen einem \textit{ServiceLocator}~\cite[S. 315-340]{ACM03}, allerdings findet der Zugriff auf die beteiligten Ressourcen über klar definierte Schnittstellen statt: Beispielsweise registriert ein Spiel über die der \textit{Soft Architecture}-Schicht den \texttt{GameObjectPoolManager} über

\begin{verbatim}
    gameWorld.registerManager<GameObjectPoolManager>();
\end{verbatim}

und greift später über

\begin{verbatim}
    auto& manager = gameWorld.manager<GameObjectPoolManager>();
\end{verbatim}

oder wahlweise


\begin{verbatim}
    auto* manager = gameWorld.tryManager<GameObjectPoolManager>();
\end{verbatim}

darauf zu.
Auch hier wird der Zugriff bzw. das Hinzufügen über \textit{Concepts} abgesichert (\texttt{IsManagerLike}).\\


Da helios auf eine Typindentifizierung über Vererbung verzichtet, muss der Klassenname der angefragten Ressource bekannt sein, womit gleichzeitig sichergestellt wird, dass von den entsprechenden Typen nur eine Instanz zur Verfügung steht.

\textit{Views}, die einen Zugriff auf eine Sammlung von Objektinstanzen ermöglichen, die unter einem bestimmten \texttt{EngineRoleTag} spezifiziert wurden, sind über semantische eindeutige und plurifizierende Methoden erreichbt, wie \texttt{managers()} und \texttt{commandHandlers()}.
Diese liefern einen \texttt{std::span} auf die über die intern über das \textit{Type Erasure}-Idiom entworfene Instanzen zurückliefern.
Die Implementierung in helios ist hier weitgehend angelehnt an~\cite{TypeErasure}.

Die GameWorld bietet außerdem Funktionen zur Erstellung von \texttt{GameObject}s an sowie die Möglichkeit, einen gefilterten View auf solche Entitäten anzufragen, die über eine bestimmte Konstellation von Komponenten verfügen, was für die in Abschnitt~\label{sec:systems} beschriebenen \textit{Systeme} relevant ist.

Als weitere Aufgabe verwaltet die \texttt{Session}.
Hierbei handelt es sich um die semantische Definition eines \texttt{GameObject}s, dem während der Laufzeit des Spiels die Rolle zukommt, auf solche Komponenten Zugriff zu erlauben, die von besondeer Bedeutung für ein Spiel sind, um zentral den Zugriff auf diese Komponenten zu erlauben; hierzu gehören unter anderem \texttt{GameState}, \texttt{MatchState} und die aktiven \texttt{ViewportId}s.